\documentclass[../../main.tex]{subfiles}
\begin{document}

2つまたは3つのベクトルの加法について，次の法則が成立する．
\begin{align*}
  \vec{A}+\vec{B}=\vec{B}+\vec{A}, \hspace{1\zw}
  (\vec{A}+\vec{B})+\vec{C} =
  \vec{A}+(\vec{B}+\vec{C})
\end{align*}
いま，$n$個のベクトルを$\vec{A}_1,\vec{A}_2,\cdots,\vec{A}_n$とし，
その順序を任意にかえたものを，$\vec{B}_1,\vec{B}_2,\cdots,\vec{B}_n$とする．
上の2つの法則だけをつかって，
\begin{align*}
  \vec{A}_1+\vec{A}_2+\cdots+\vec{A}_n =
  \vec{B}_1+\vec{B}_2+\cdots+\vec{B}_n
\end{align*}
がなりたつことを数学的帰納法を用いて示せ．
なお，
たとえば4つのベクトル$\vec{A}$，$\vec{B}$，$\vec{C}$，$\vec{D}$について，
その和$\vec{A}+\vec{B}+\vec{C}+\vec{D}$は
$ \{ (\vec{A}+\vec{B})+\vec{C} \} +\vec{D}$を意味するものとし，
一般の場合も同様とする．

\end{document}
